% ============================================================
\section{顶刊前沿：近年 AER/JPE 怎样用效用函数（深化版）}
% ============================================================

本节基于 2021--2026 年 AER、JPE 及相关顶刊的公开文献（均经联网核实），
按 \textbf{文献 → 模型（公式 + 假设）→ 新质 → 可借鉴处} 拆解。
重点不在“他们用了哪个名字的形式”，而在\textbf{效用函数的自变量、
参数化与识别策略}如何被重新设计。

\subsection{随机系数 Logit 的偏好面板化：D\"opper--MacKay--Miller--Stiebale（JPE 2025）}\label{subsec:jpe-markups}

\notehead{模型}
个体 $i$ 在市场 $t$（品类 $\times$ 年）购买产品 $j$ 的间接效用：
\[
u_{ijt}=\underbrace{x_{jt}\beta-\alpha p_{jt}+\xi_{jt}}_{\delta_{jt}\ \text{(均值效用)}}
+\underbrace{\sum_k x_{jkt}\,\sigma_k\,\nu_{ik}}_{\mu_{ijt}\ \text{(个体异质)}}
+\varepsilon_{ijt},
\]
其中 $\nu_{ik}\sim N(0,1)$ 为个体口味冲击，$\sigma_k$ 度量偏好的离散程度，
$\varepsilon_{ijt}$ 为 I 型极值扰动（§\ref{subsec:rum}）。

\notehead{关键假设}
(i) $\mathbb{E}[\xi_{jt}\mid Z_{jt}]=0$：工具变量（成本、竞品特征）与不可观测质量正交；
(ii) 企业按 Nash--Bertrand 定价：价格加成 $=$ 需求弹性倒数（一阶条件），
从而\textbf{边际成本可由需求参数反推}；
(iii) 偏好在品类内随年份变化——每个品类 $\times$ 年\textbf{单独}估计
$(\alpha_t,\sigma_t,\delta_{jt})$。

\notehead{新质}
传统用法把 $(\alpha,\sigma)$ 当跨期结构常数；本文得到\textbf{偏好面板}：
2006--2019 年 $\alpha_t$（价格敏感度）系统性下降，
对应消费者剩余上移与企业加成率平均上升约 30\%。
偏好参数从“待估常数”变成“待解释的时变对象”——
这是需求估计范式的转向。

\notehead{可借鉴}
任何长期政策评估（养老金、医疗市场）若样本跨十年以上，
“偏好稳定”假设应先被检验而非默认；
检验成本不高：分时段重估并比较参数即可。

\subsection{福利分析对函数形式的稳健性：Kang--Vasserman（AER 2025）}

\notehead{模型}
不指定需求函数形式。设真实需求 $q(p)$ 未知，
分析者假设某形式 $\hat q(p)$ 并据此算福利 $\hat W$。
定义\textbf{稳健半径}：在梯度与曲率允许偏差
$\|\nabla q-\nabla\hat q\|\le\Delta_1$、$\|\nabla^2 q-\nabla^2\hat q\|\le\Delta_2$
内，使福利结论反转所需的\textbf{最小} $(\Delta_1,\Delta_2)$。
用信息设计工具把最小化问题化为可计算规划。

\notehead{新质}
把“换两组参数重算”的朴素稳健性升级为\textbf{形式无关的度量}：
回答的不是“参数错了怎么办”，而是“形式错了多远结论才翻”。

\notehead{可借鉴}
政策论文的福利部分可附一行稳健半径：
“需求曲率需偏离假设形式 X\% 以上，本结论才反转”——
比传统敏感性表格信息密度高一个量级（§4.8 第 9 步已纳入）。

\subsection{非位似 CES 与结构转型：Comin--Lashkari--Mestieri（ECMA 2021 及 2025 后续）}

\notehead{模型}
偏好由隐式加法可分形式定义（§4.2.3）：
$\sum_i\Omega_i(c_iU^{\varepsilon_i})^{(\sigma-1)/\sigma}=1$。
关键可估对象：相对恩格尔曲线斜率 $\varepsilon_i-\varepsilon_j$，
用美国家庭、印度家庭及 39 国历史面板分别估计。

\notehead{关键假设}
(i) 部门收入弹性 $\varepsilon_i$ 为常数（恩格尔斜率不衰减）；
(ii) 价格效应（$\sigma$）与收入效应（$\varepsilon_i$）分离进入。

\notehead{新质与结论}
收入效应解释了\textbf{国家内部}结构转型的大部分；
位似偏好（$\varepsilon_i\equiv1$）在数据上被明确拒绝。
2025 年后续工作把同一偏好嵌入创新模型：
收入弹性高的部门（服务）吸引更多研发——
需求结构\textbf{引导}技术进步方向。

\notehead{可借鉴}
研究中国消费升级/产业升级时，位似 CES 的“份额恒定”
假设直接排除核心机制；非位似 CES 是当前可估计的标准件。

\subsection{选择括定的显示偏好检验：Ellis--Freeman（AER 2024）}

\notehead{模型}
观察个体在一系列\textbf{同时呈现的}预算集上的选择。
若个体宽括定：所有预算的联合约束下最大化单一效用；
若窄括定：每个预算独立最大化各自的（可能不同的）效用。
文章给出两组选择数据上的公理（弱/强括定的一致性条件），
无需指定效用函数参数形式即可判定括定类型。

\notehead{新质}
把“心理账户”从实验设定变量变为\textbf{可从常规选择数据识别}的对象；
方法完全非参数（与本课程第二章显示偏好路线同源）。

\notehead{可借鉴}
养老金、医疗共付、教育储蓄等“天然多账户”场景，
建模前应先问括定方式；宽/窄括定下同一政策的福利评估可以反向。

\subsection{随机效用的可观测性边界：Kono--Saito--Sandroni（AER 2026）}

\notehead{模型}
标准 RUM（§\ref{subsec:rum}）假设所有备选的选择频率可观测；
现实中“外部选项”（未购买、未申报）的频率常缺失。
文章推导：当部分方案频率不可观测时，
刻画 RUM 可理性化的\textbf{有限、非冗余不等式组}。
核心发现：把不可观测方案笼统聚合为一个“外部选项”
会丢失 RUM 的可检验内容——\textbf{即使数据与聚合后的 RUM 一致，
也可能不存在任何与原始数据一致的随机效用模型}。

\notehead{可借鉴}
用 Logit 族估计前，先检查数据的外部选项处理方式；
外部选项的定义（包含什么、频率多少）本身是识别假设。

\subsection{序数结果上的社会偏好：Ambuehl--Bernheim（AER 2026）}

\notehead{模型}
替他人决策时，决策者知道他人对方案的\textbf{排序}但不知强度。
实验发现：多数人对“避免他人的最差结果”赋的权重，
高于“促成他人的最好结果”——排序信息被\textbf{位置性}使用
（类似 §4.6.2 最大最小倾向，但建立在序数信息上），
且该原则跨国稳定。

\notehead{可借鉴}
代议决策（家长为子女、受托人为受益人、政府为公民）
的福利函数不该是基数值的简单加总；
“防止最差”权重项（如 $\min$ 型或位置加权）更贴近真实社会偏好。

\subsection{模糊偏好的更新：Sadowski--Sarver（AER: Insights 2026）}

\notehead{模型与新质}
多先验/非期望效用（§\ref{subsec:meu}）在动态环境下有内在张力：
更新规则不能同时满足\textbf{动态一致}与\textbf{结果主义}
（只按事后偏好行事）。文章用演化方法论证：
动态一致是演化最优的硬要求，而结果主义的系统性违反反而最优——
为“事先承诺 vs 事后反悔”提供了演化基础。

\notehead{可借鉴}
模糊模型用于政策评估时，必须显式声明更新规则；
不同的更新规则下，同一政策的福利排序可以不同。

\subsection{异质个体的福利评估与有限注意：JPE 2025 两例}

\begin{itemize}
  \item \textbf{D\'avila \& Schaab}（JPE 2025）：
        个体在偏好与约束上多维异质时，福利变化分解为
        个体层面边际福利权重的加总——
        纠正“代表性消费者”对再分配效应的系统性遗漏；
  \item \textbf{Angeletos \& Sastry}（JPE 2025）：
        需求函数的斜率依赖\textbf{注意分配}：
        同一效用函数下，价格是否被注意决定需求弹性，
        名义刚性获得偏好以外的微观基础。
\end{itemize}

\subsection{中国问题中的效用函数：延迟退休与生育（详见 \S 4.6）}

\begin{center}
\footnotesize
\setlength{\tabcolsep}{4pt}
\begin{tabular}{p{1.9cm}p{3.0cm}p{4.6cm}p{4.6cm}}
\toprule
文献 & 模型 & 效用形式 & 关键结论 \\
\midrule
MDPI（2024） & 两期 OLG + 内生生育 & $\ln c_1+\beta\ln c_2+\phi\ln n+\eta\ln d$ & 延迟退休、缴费率与家庭赡养转移负相关 \\
SSRN（2024） & 多期 OLG + 祖辈照护 & 内生生育 + 隔代照料时间 & 延迟退休 5 年或使生育率降逾 25\% \\
IMF（2026） & 80 代 OLG、城乡分割 & 生命周期 CRRA 型 & 延迟退休 3 年：GDP +5.6\%（2050），养老金支出占比 15.3$\to$11.9\% \\
系统工程（2024） & OLG 内生生育与退休 & 生育补贴入预算约束 & 补贴促生育且促延迟退休；养老金呈倒 U \\
\bottomrule
\end{tabular}
\end{center}
共同特征：全部使用对数/CRRA 可加效用以换取闭式解与可校准性；
差异在\textbf{效用的自变量}——是否包含子女数量 $n$、父母消费 $d$、
祖辈照料时间。§4.9 完整复现此类模型的推导与计算。

\subsection{前沿速查表（2021--2026）}

\begin{center}
\small
\begin{tabular}{lll}
\toprule
方向 & 传统写法 & 前沿改写 \\
\midrule
福利稳健性 & 指定形式 $\to$ 算福利 & 形式无关的稳健半径（AER 2025） \\
偏好参数 & 结构常数 & 时变偏好面板（JPE 2025） \\
收入效应 & 位似（份额恒定） & 非位似、恩格尔斜率恒定（ECMA 2021） \\
选择架构 & 全合并、全理性 & 括定可检验（AER 2024） \\
数据基础 & 外部选项随手聚合 & 可观测性边界（AER 2026） \\
社会偏好 & 基数加总 & 序数位置加权（AER 2026） \\
模糊与动态 & 单一先验 & 多先验 + 更新规则声明（AERI 2026） \\
个体性 & 代表性消费者 & 多维异质 + 注意内生（JPE 2025） \\
\bottomrule
\end{tabular}
\end{center}
